\scanpage{93}
$I'$ sur $\widehat X'$, on peut aussi, plus généralement, se donner une partie $I'$ (non pas de $I$, mais de) $\widetilde J\amalg I$, qu'on \unsure{retiendra} comme ``pts marqués''.

Quand on s'est ainsi donné (pour un \deleted{\unsure{MD}}m.graphe $\underline G$ donné) une partie $A'$ de $A(G)$, et une partie $I'$ de $I\amalg\widetilde J$ (où $\widetilde J=\widetilde A(G)-\widetilde A'\deleted{\unsure{)}}$), alors pour la courbe spéciale $G$-épinglée, on déduit une courbe spéciale, qui sera épinglée par le m.graphe $\underline G'$ déduit de $\underline G,A',I'$ en prenant le sous-graphe \inserted{$G'$ de $G$} ayant \unsure{$=$|les} $S$ sommets et ayant $A'$ comme ens. d'arêtes, avec comme ``ens de pts marqués'' $I'\subset I\amalg\widetilde J$, qui s'envoie dans $S$ \deleted{\unsure{si}} par
\[
(120)\qquad I'\longrightarrow S\quad\text{induit par }I\longrightarrow S\text{ et }\widetilde A(G)\xrightarrow{\text{origine}}S,
\]
l'application genre pour $\underline G'$ est celle pour $\underline G$.
\margin{92}
